\notechapter{N-20230127}{Factorization of \texorpdfstring{$\sum_{i=0}^k x^{in}$}{sum x^{in}}: Roots of Unity, False Generalization, and Cyclotomic Structure}

\section{The original problem and the first successful method}

The paper begins with the factorization of
\[
  x^{2n}+x^n+1
\]
by
\[
  x^2+x+1.
\]
Let $\omega$ be a primitive cubic root of unity.  Then
\[
  \omega^2+\omega+1=0,\qquad \omega^3=1,\qquad \omega\ne1.
\]
The polynomial $x^2+x+1$ divides $x^{2n}+x^n+1$ iff both $\omega$ and $\omega^2$ are roots.  Substituting $x=\omega$ gives
\[
  \omega^{2n}+\omega^n+1.
\]
This is zero exactly when $\omega^n$ is a nontrivial cubic root of unity.  Therefore the divisibility holds iff
\[
  n\not\equiv0\pmod3.
\]
This proof is short because roots of unity convert a divisibility question into an exponent congruence question.

\section{The direct factorization attempt}

The note then tries to prove the same statement by direct algebra.  Since
\[
  x^{2n}+x^n+1=\frac{x^{3n}-1}{x^n-1},
\]
one can try to factor numerator and denominator using difference of squares or cubes.  For example, when $n=2$,
\[
  x^4+x^2+1=(x^2+x+1)(x^2-x+1).
\]
When $n=4$, one obtains several quadratic-looking factors, including again $x^2+x+1$.

But when $n$ changes parity or becomes divisible by $3$, the algebraic route splits into cases.  The original article carefully lists cases such as $n=3m\pm1$ and tracks which geometric-series quotients appear.  The lesson is that direct factorization can prove the theorem, but it obscures the simple reason: the exponent map $z\mapsto z^n$ acts on the three roots of unity.

\section{The attempted generalization}

The natural next question is whether
\[
  1+x^n+x^{2n}+\cdots+x^{kn}
\]
is divisible by
\[
  1+x+x^2+\cdots+x^k
\]
whenever $n$ is not a multiple of $k+1$.  The author first guesses yes, then uses roots of unity to test the guess.

Let $\xi$ be a nontrivial $(k+1)$-st root of unity.  We need
\[
  1+\xi^n+\xi^{2n}+\cdots+\xi^{kn}=0
\]
for every such $\xi$.  The geometric sum vanishes if $\xi^n\ne1$.  Thus the required condition is that no nontrivial $(k+1)$-st root becomes $1$ after raising to the $n$th power.  This is equivalent to
\[
  \gcd(n,k+1)=1.
\]
So the first guess $n\not\equiv0\pmod{k+1}$ is too weak.  For example, $n$ can be nonzero modulo $k+1$ but still share a nontrivial common divisor with $k+1$.

\section{The computational experiment}

The note then turns to MATLAB.  It computes roots of
\[
  x^{kn}+x^{(k-1)n}+\cdots+x^n+1=0
\]
for different values of $n$ and plots them in the complex plane, sometimes using a third coordinate to record $n$.  This is not just decoration.  The picture reveals periodicity in the arguments of the roots.

The observed pattern is: the roots lie on the unit circle, their arguments are fractions of $2\pi$, and as $n$ varies the same angular patterns repeat.  For $k=2$, the pattern has two angles per block; for $k=3$, three angles per block; in general, the $(k+1)$-st-root structure controls the blocks.

This experimental section is important because it records a genuine mathematical process: a false conjecture is tested, numerical pictures suggest the right invariant, and then the invariant is expressed exactly by roots of unity.

\section{Explicit root formula}

The equation
\[
  1+x^n+x^{2n}+\cdots+x^{kn}=0
\]
can be rewritten as
\[
  \frac{x^{n(k+1)}-1}{x^n-1}=0.
\]
Thus the roots are those $n(k+1)$-st roots of unity whose $n$th power is not $1$.  Equivalently,
\[
  x=\exp\left(\frac{2\pi i}{n(k+1)}(r(k+1)+m)\right),
\]
where $r=0,\ldots,n-1$ and $m=1,\ldots,k$.  This gives exactly $kn$ roots, as expected from the degree.

The original note writes these roots as products of exponentials:
\[
  e^{2\pi i r/n}e^{2\pi i m/(n(k+1))}.
\]
Substituting into the polynomial gives a finite geometric sum whose ratio is a nontrivial $(k+1)$-st root of unity, hence the sum is zero.

\section{Common real factors}

The final theorem concerns common real-coefficient factors among these polynomials as $n$ varies.  Since complex roots come in conjugate pairs, a set of roots stable under conjugation gives a real polynomial factor.  When $d=\gcd(n,k+1)>1$, some of the angular positions line up across different $n$, producing common roots and hence common real factors.

This is the cyclotomic viewpoint in disguise.  The polynomial
\[
  x^m-1
\]
factors as
\[
  x^m-1=\prod_{d\mid m}\Phi_d(x),
\]
and the geometric sum
\[
  1+x+\cdots+x^k=\frac{x^{k+1}-1}{x-1}
\]
contains the cyclotomic factors $\Phi_d(x)$ for divisors $d>1$ of $k+1$.  Replacing $x$ by $x^n$ pulls these cyclotomic conditions back along the power map.

\section{The organizing idea}

The value of this note is the exploration, not only the final condition.  It begins with a correct theorem, tries a direct factorization proof, guesses a generalization, finds that the guess is false, uses computation to see the roots, and finally explains the pattern by complex exponentials.  That is a miniature version of mathematical research: algebraic manipulation, conjecture, counterexample, visualization, and structural explanation.

\section{Geometric sums and roots of unity}
The polynomial
\[
1+x+\cdots+x^m=\frac{x^{m+1}-1}{x-1}
\]
has roots among roots of unity except \(1\).  Factorization questions are therefore questions about which roots of unity two polynomials share.
\section{Cyclotomic organization}
The factorization
\[
x^n-1=\prod_{d\mid n}\Phi_d(x)
\]
separates roots by exact order.  Common factors of expressions built from geometric sums can be read by comparing divisibility relations among orders.
\section{Explicit roots and common factors}
If roots are written as
\[
e^{2\pi i k/n},
\]
then common real factors come from conjugate pairs.  A real polynomial factor corresponds to a set of roots closed under complex conjugation.
\section{Cyclotomic structure behind geometric sums}
The computational experiments point toward the same structure: geometric-sum factorizations are controlled by roots of unity, cyclotomic polynomials, orders, and conjugation symmetry.

\section{Cyclotomic polynomials}

The roots of $x^n-1$ are the $n$-th roots of unity.  The primitive $n$-th roots are those with exact order $n$.  The cyclotomic polynomial is
\[
  \Phi_n(x)=\prod_{\substack{1\le k\le n\\ \gcd(k,n)=1}}(x-e^{2\pi i k/n}).
\]
It has degree $\varphi(n)$ and integer coefficients.

The factorization is
\[
  x^n-1=\prod_{d\mid n}\Phi_d(x).
\]
This identity is the structural version of many elementary root-of-unity factorization exercises.

\section{Galois action on roots of unity}

The field $\mathbb Q(\zeta_n)$ generated by a primitive root of unity $\zeta_n$ is the $n$-th cyclotomic field.  Its Galois group is
\[
  \operatorname{Gal}(\mathbb Q(\zeta_n)/\mathbb Q)\cong(\mathbb Z/n\mathbb Z)^\times.
\]
An element $a\in(\mathbb Z/n\mathbb Z)^\times$ acts by
\[
  \zeta_n\mapsto \zeta_n^a.
\]
Thus the same reduced residue system that appears in elementary number theory controls the symmetries of roots of unity.

\section{Resultants and common roots}

Two polynomials have a common root exactly when their resultant is zero.  For example, studying when a cyclotomic polynomial divides another expression can be reframed as asking which roots are common.

The elementary method says: substitute a primitive root $\omega$ and check whether the expression vanishes.  The algebraic version says: divisibility is detected by vanishing on all Galois conjugate roots, or by a resultant condition.

\section{Finite fields and character sums}

Roots of unity also appear in finite fields.  The multiplicative group $\mathbb F_q^\times$ is cyclic of order $q-1$.  Therefore roots of $x^n=1$ in $\mathbb F_q$ are controlled by $\gcd(n,q-1)$.

Character sums such as
\[
  \sum_{x\in\mathbb F_q}\chi(x)\psi(f(x))
\]
are research-level descendants of elementary root-of-unity sums.  They connect finite fields, Fourier analysis on finite groups, and algebraic geometry.

\section{Returning to the original experiments}

The note's factorization experiments with roots on the complex plane are the concrete beginning of cyclotomic theory.  The advanced bridge is precise: roots of unity form finite cyclic groups, cyclotomic polynomials package primitive roots, and Galois groups act by exponentiation modulo $n$.

\section{Primitive roots and cyclotomic factor tests}

To decide whether a polynomial expression is divisible by $\Phi_n(x)$, it is enough to substitute a primitive $n$-th root of unity $\zeta$ and check whether the expression vanishes.  Because $\Phi_n$ is the minimal polynomial of $\zeta$ over $\mathbb Q$, vanishing at one primitive root forces vanishing at all Galois conjugates.

This explains why root-of-unity substitution is not just a clever complex-number trick.  It is a minimal-polynomial test.

\section{Cyclotomic identities through Möbius inversion}

From
\[
  x^n-1=\prod_{d\mid n}\Phi_d(x)
\]
one obtains
\[
  \Phi_n(x)=\prod_{d\mid n}(x^d-1)^{\mu(n/d)}.
\]
This is Möbius inversion on the divisor lattice.  The alternating exponents in this formula are the same inclusion--exclusion pattern seen in counting and arithmetic functions.

Thus the note's factorization experiments sit at an intersection of complex roots, integer polynomials, Galois symmetry, and Möbius inversion.
